% !TeX root = ../main.tex

% 中文摘要
\begin{abstract}
含参函数的“任意性”与“存在性”问题，即函数“恒成立”与“能成立”问题，是高考数学中的高频考点，也是高中数学教师教学与学生学习的难点问题。此类问题综合性强、解法多样，对学生的逻辑思维和数学素养要求较高。本文以优化教学策略、提升教学质量为目标，通过系统研究，旨在为师生提供更系统的解题思路和教学方法。

本文从三个维度对问题进行分类整理：一是单函数与双函数的区分，二是单变量与双变量的讨论，三是任意性与存在性的量词组合。通过分类归纳，总结出典型题型及其解法特征。并研究采用“问题驱动、直观感知、逻辑推理、类比迁移”的专题教学策略，旨在帮助学生构建清晰的知识网络，提升分析问题和解决问题的能力。

目前针对此类问题的教学研究较为零散，缺乏系统的题型梳理和统一框架。本文的创新点在于：提出了一种全面、系统的题型分类方法，弥补了现有研究的不足；采用一道核心例题的渐进变式串联所有题型，首次构建了此类问题的统一框架，揭示了不同题型间的内在联系；设计了一套可视化教学方案，将数形结合融入教学，帮助学生理解题意，培养其逻辑推理和综合应用能力，也为高中数学教师提供了教学参考。
\end{abstract}



% 英文摘要
\begin{abstract*}
  The ``arbitrariness'' and ``existence'' issues of parameter-containing functions, namely the ``always true'' and ``can be true'' problems of functions, are high-frequency examination points in the college entrance examination of mathematics and also difficult points for high school mathematics teachers in teaching and students in learning. Such problems are highly comprehensive and have diverse solutions, which require a high level of logical thinking and mathematical literacy from students. This article aims to optimize teaching strategies and improve teaching quality. Through systematic research, it is intended to provide more systematic problem-solving ideas and teaching methods for both teachers and students.

  This paper classifies and organizes the problems from three dimensions: firstly, the distinction between single-function and double-function; secondly, the discussion on single-variable and double-variable; thirdly, the combination of quantifiers of arbitrariness and existence. Through classification and induction, typical problem types and their solution characteristics are summarized. Moreover, it studies the application of the thematic teaching strategy of ``problem-driven, intuitive perception, logical reasoning, and analogy transfer'', aiming to help students construct a clear knowledge network and enhance their abilities in analyzing and solving problems.

  At present, the teaching research on such issues is rather scattered and lacks systematic classification of question types and a unified framework. The innovation points of this paper lie in: proposing a comprehensive and systematic classification method for question types, which makes up for the deficiencies of existing research; using a core example with progressive variations to connect all question types, and for the first time constructing a unified framework for such problems, revealing the intrinsic connections among different question types; designing a set of visual teaching schemes, which assist students in understanding the problem context through the combination of numbers and shapes, cultivating their logical reasoning and comprehensive application abilities, and also providing teaching references for high school mathematics teachers.
\end{abstract*}